\section*{\texorpdfstring{\S\CurrentExerciseGroup}{第{\CurrentExerciseGroup}节}}
\phantomsection
\addcontentsline{toc}{section}{第{\CurrentExerciseGroup}节习题}

\begin{enumerate}
\def\labelenumi{\arabic{enumi})}
\item
  设 \(\mathrm{T}, \mathrm{T}'\) 为两个局部紧空间，\(\mu\) 为 \(\mathrm{T}\) 上的正测度，\(\mu'\) 为 \(\mathrm{T}'\) 上的正测度。证明：若测度 \(\mu'\) 有界，则投影 \(\mathrm{pr}_1\) 把 \(\mathrm{T} \times \mathrm{T}'\) 映到 \(\mathrm{T}\) 上，它是 \((\mu \otimes \mu')\) -逆紧映射，且 \(\mathrm{pr}_1(\mu \otimes \mu') = a \cdot \mu\)，其中 \(a = \mu'(\mathrm{T}')\)。由此给出一个 \((\mu \otimes \mu')\) -可忽略子集（含于 \(\mathrm{T} \times \mathrm{T}'\)）的例子，使它在 \(\mathrm{T}\) 上的投影不是 \(\mu\) -可测的。
\item
  设 T 为 R 中区间 \([0,1]\)，赋以由 R 的拓扑诱导的拓扑，设 \(T'\) 为 R 中区间 \([0,1]\)，赋以离散拓扑；设 \(\mu\) 为 T 上的 Lebesgue 测度，\(\mu'\) 为 \(T'\) 上由 \(T'\) 的每点质量 +1 定义的离散测度，并设 \(\nu = \mu \otimes \mu'\) 为 \(X = T \times T'\) 上的乘积测度。
\end{enumerate}

\begin{enumerate}
\def\labelenumi{\alph{enumi})}
\item
  对每个 \(t \in \mathrm{T}\)，集合 \(\{t\}\) 是 \(\mu\) -可忽略的，但集合 \(\{t\} \times \mathrm{T}'\) 不是 \(\nu\) -可忽略的。
\item
  “对角线”\(\Delta\)（\(\mathrm{T} \times \mathrm{T}'\) 中点 \((t, t)\) 所成的集，其中 \(t\) 跑遍 \([0, 1]\)）是 \(\mathrm{T} \times \mathrm{T}'\) 中的闭集，局部 \(\nu\) -可忽略但非 \(\nu\) -可忽略，且 \(\mu'(\Delta(t)) = 1\)（对所有 \(t \in \mathrm{T}\)）而 \(\mu(\Delta(t')) = 0\)（对所有 \(t' \in \mathrm{T}'\)）（参见 §3, 习题 3）。
\item
  在空间 \(\mathrm{X} \times \mathrm{X} = \mathrm{Y}\) 上考察乘积测度 \(\rho = \nu \times \nu\)。证明 \(\mathrm{Y}\) 中存在一个集合 \(\mathrm{A}\)，局部 \(\rho\) -可忽略，使得对每个 \(x \in \mathrm{X}\)，\(\mathrm{A}(x)\) 与 \(\mathrm{A}^{-1}(x)\) 都是 \(\nu\) -可测的且非局部可忽略的。（若 \(x = (t, t')\)，其中 \(t \in \mathrm{T}\)，\(t' \in \mathrm{T}'\)，取 A 使 A(x) 包含所有点 \((s, t) \in X\)，其中 \(0 \leqslant s \leqslant 1\)。）
\end{enumerate}

\begin{enumerate}
\def\labelenumi{\arabic{enumi})}
\setcounter{enumi}{2}
\tightlist
\item
  设 T 为局部紧空间，\(\mu\) 为 T 上的正测度，f 为定义在 T 上的数值函数 \(\geqslant 0\)。在乘积空间 \(T \times R\) 中，用 \(D_{f}\) 表示点 \((t, x)\) 所成之集，其中 \(0 \leqslant x \leqslant f(t)\)；另一方面，设 \(\nu\) 为 R 上的 Lebesgue 测度。
\end{enumerate}

\begin{enumerate}
\def\labelenumi{\alph{enumi})}
\item
  证明：为使 \(f\) 是 \(\mu\) -可测的，必须且只需 \(D_f\) 是关于乘积测度 \(\lambda = \mu \otimes \nu\) 的可测集。（为证条件必要，证明：若 \(f\) 是 \(\mu\) -可测的，则对每个紧子集 \(K\)（含于 \(T\)），\(D_f \cap (K \times R_+)\) 是一个 \(\lambda\) -可忽略集与形如 \(A \times I\) 的集所成的一个可数族之并，其中 \(A\) 是 \(\mu\) -可测的，\(I\) 是 \(R_+\) 的一个区间。为证条件充分，证明：若该条件满足，则对每个紧子集 \(K\)（含于 \(T\)），存在稠密集 \(H \subset R_+\)，使得对每个 \(\alpha \in H\)，\(K \cap f^{-1}([\alpha, +\infty])\) 是 \(\mu\) -可测的，这里要用命题 7 的推论 2。）
\item
  证明：为使 \(f\) 是 \(\mu\) -可积的，必须且只需 \(\mathrm{D}_f\) 是 \(\lambda\) -可积的，且此时有 \(\lambda(\mathrm{D}_f) = \int f d\mu\)。此外，用 \(g\) 表示 \(\mathbf{R}_+\) 上由 \(g(t) = \mu\left( f([t, +\infty]) \right)\) 定义的递减数值函数（可能等于 \(+\infty\)，当 \(t = 0\) 时（Ch. IV, §5, 习题 29)），则有 \(\int f d\mu = \int_0^{+\infty} g(t) dt\)。
\end{enumerate}

\begin{enumerate}
\def\labelenumi{\arabic{enumi})}
\setcounter{enumi}{3}
\tightlist
\item
  \begin{enumerate}
  \def\labelenumii{\alph{enumii})}
  \tightlist
  \item
    设 \(T, T'\) 为两个无穷远处可数的局部紧空间，\(\mu\) 为 T 上的正测度，\(\mu'\) 为 \(T'\) 上的正测度。证明：对每个 \((\mu \otimes \mu')\) -可测映射 f（把 \(T \times T'\) 映入 R 中），函数 \(F : t \mapsto \int^{*} |f(t, t')| d\mu'(t')\) 是 \(\mu\) -可测的（把 \(T'\) 看作递增紧集序列之并）。
  \end{enumerate}
\end{enumerate}

\begin{enumerate}
\def\labelenumi{\alph{enumi})}
\setcounter{enumi}{1}
\tightlist
\item
  再假设，对几乎每个 \(t \in T\)（关于 \(\mu\)）函数 \(f(t, \cdot)\) 是 \(\mu'\) -可积的，且对几乎每个 \(t' \in T'\)（关于 \(\mu'\)）函数 \(f(\cdot, t')\) 是 \(\mu\) -可积的；最后，假设函数（关于 \(\mu\) 几乎处处有定义）
\end{enumerate}

\(t \mapsto \int f(t, t') d\mu'(t')\) 是 \(\mu\) -可积的。证明函数（关于 \(\mu'\) 几乎处处有定义）\(t' \mapsto \int f(t, t') d\mu(t)\) 此时是 \(\mu'\) -可积的，且

\[
\int d \mu^ {\prime} (t ^ {\prime}) \int f (t, t ^ {\prime}) d \mu (t) = \int d \mu (t) \int f (t, t ^ {\prime}) d \mu^ {\prime} (t ^ {\prime}).
\]

（由 \(a\)），存在 \(\mathbf{T}\) 的一个分割，由一个 \(\mu\) -可忽略集 \(\mathbf{N}\) 与一列紧集 \((\mathrm{K}_n)\) 组成，使函数 \(f\varphi_{\mathrm{K}_n\times \mathrm{T}'}\) 中每一个都是 \((\mu \otimes \mu')\) -可积的。令

\[
g _ {n} (t ^ {\prime}) = \sum_ {j = 1} ^ {n} \int_ {K _ {j}} f (t, t ^ {\prime}) d \mu (t)
\]

并应用 §5 习题 20。）

\begin{enumerate}
\def\labelenumi{\alph{enumi})}
\setcounter{enumi}{2}
\tightlist
\item
  设 \(\mathbf{T}\) 为 \([0,1]\)（\(\mathbf{R}\) 中的区间），\(\mu\) 为 \(\mathbf{T}\) 上的 Lebesgue 测度。对每个整数 \(n > 0\)，令 \(\mathrm{A}_n^\prime = \left[\frac{1}{2^n}, \frac{3}{2^{n+1}}\right]\)，\(\mathrm{A}_n^{\prime\prime} = \left[\frac{3}{2^{n+1}}, \frac{1}{2^{n-1}}\right]\)，并在乘积空间 \(\mathbf{T} \times \mathbf{T}\) 中令 \(\mathrm{B}_n^\prime = \mathrm{A}_n^\prime \times \mathrm{A}_n^\prime\)，\(\mathrm{B}_n^{\prime\prime} = \mathrm{A}_n^{\prime\prime} \times \mathrm{A}_n^{\prime\prime}\)，\(\mathrm{C}_n^\prime = \mathrm{A}_n^\prime \times \mathrm{A}_n^{\prime\prime}\)，\(\mathrm{C}_n^{\prime\prime} = \mathrm{A}_n^{\prime\prime} \times \mathrm{A}_n^\prime\)。令 \(f(t,t') = 4^{n+1}\) 于 \(\mathrm{B}_n^\prime\) 与 \(\mathrm{B}_n^{\prime\prime}\) 中，令 \(f(t,t') = -4^{n+1}\) 于 \(\mathrm{C}_n^\prime\) 与 \(\mathrm{C}_n^{\prime\prime}\) 中（对每个整数 \(n > 0\)），并令 \(f(t,t') = 0\) 于 \(\mathbf{T} \times \mathbf{T}\) 的其余点处。证明两个积分 \(\int d\mu(t') \int f(t,t') d\mu(t)\) 与 \(\int d\mu(t) \int f(t,t') d\mu(t')\) 都有定义且相等，但函数 \(f\) 关于测度 \(\mu \otimes \mu\) 不可积。
\end{enumerate}

¶ 5) 设 T, T′ 为两个局部紧空间，μ 为 T 上的正测度，μ′ 为 T′ 上的正测度；假设 T 的每个紧子集都可度量化。设 f 为 T × T′ 到可度量化空间 G 中的映射，满足：1° 对每个 t ∈ T，映射 t′ → f(t, t′) 是 μ′-可测的；2° 对每个 t′ ∈ T′，映射 t → f(t, t′) 是连续的。证明在这些条件下，f 是 (μ ⊗ μ′)-可测的。（指出（利用 Egoroff 定理）：对 T 的每个紧子集 K，f 在 K × T′ 上的限制是一列 (μ ⊗ μ′)-可测函数的极限。）

¶ 6) a) 设 T 为局部紧空间，ν 为 T 上的正测度，I 为 R 的区间，μ 为 I 上的正测度。设 A 为乘积 I × T 的一个子集，满足：1° 对每个 x ∈ I，A 在 x 处的截面 A(x) ⊂ T 是 ν-可测的；2° 关系 x ≤ y 蕴含 A(x) ⊂ A(y)。证明 A 是 (μ ⊗ ν)-可测的。（化归于 I 与 T 紧且 μ 弥散的情形；考察 I 的点的递增序列 (xi)0≤i≤n，x0 与 xn 为 I 的端点，使对 0 ≤ i ≤ n-1 有 μ({[}xi, xi+1{]}) ≤ ε，然后作 I × T 的两个子集 B, C，关于 λ = μ ⊗ ν 可测，使 B ⊂ A ⊂ C 且 λ(C - B) ≤ εν(t)。）

\begin{enumerate}
\def\labelenumi{\alph{enumi})}
\setcounter{enumi}{1}
\item
  由 a) 推出：若 \(f\) 是定义在 \(\mathrm{I} \times \mathrm{T}\) 上的数值函数，满足：\(1^{\circ}\) 对每个 \(x \in \mathrm{I}\)，\(t \mapsto f(x, t)\) 是 \(\nu\) -可测的，且 \(2^{\circ}\) 对每个 \(t \in \mathrm{T}\)，\(x \mapsto f(x, t)\) 是递增的，则 \(f\) 是 \((\mu \otimes \nu)\) -可测的。
\item
  设 \(g\) 为定义在 \(I \times T\) 上的数值函数，满足：\(1^{\circ}\) 对每个 \(x \in I\)，\(t \mapsto g(x, t)\) 是 \(\nu\) -可测的；\(2^{\circ}\) 对每个 \(t \in T\)，\(x \mapsto g(x, t)\) 是单调的。由 \(b\) 推出 \(g\) 是 \((\mu \otimes \nu)\) -可测的。（设 \(T_1 \subset T\) 为 \(t \in T\) 中使映射 \(x \mapsto g(x, t)\) 递增者所成的集；证明 \(T_1\) 是 \(\nu\) -可测的。为此，对每对有理数 \((r_1, r_2)\)（属于 \(I\) 且满足 \(r_1 \leqslant r_2\)），考察 \(D_{r_1, r_2}\)，即由点 \(t \in T\) 所成的集，其中 \(g(r_1, t) \leqslant g(r_2, t)\)，并把 \(T_1\) 表为诸集 \(D_{r_1, r_2}\) 的交。）
\end{enumerate}

¶ 7) a) 设 T, T′ 为两个局部紧空间；假设 T′ 具有可数基。设 μ 为 T 上的正测度，μ′ 为 T′ 上的正测度。设 f 为定义在 T × T′ 上的数值函数 ≥0，在 T × T′ 的每个紧子集上有界，且满足：\(1^{\circ}\) 对几乎每个 \(t \in \mathrm{T}\)，函数 \(t' \mapsto f(t, t')\) 是 \(\mu'\) -可测的；\(2^{\circ}\) 对每个函数 \(h \in \mathcal{K}(\mathrm{T}')\)，函数

\[
t \mapsto \int f (t, t ^ {\prime}) h (t ^ {\prime}) d \mu^ {\prime} (t ^ {\prime}),
\]

几乎处处有定义，是 \(\mu\) -可测的。证明在这些条件下，存在一个函数 \(g\)（关于 \(\mu \otimes \mu'\) 可测），使对每个 \(t \in \mathrm{T}\)，有 \(f(t, t') = g(t, t')\)，但一个 \(\mu'\) -可忽略集 \(\mathrm{A}_t\)（依赖于 \(t\)）中的点除外。（证明：对每个函数 \(\varphi \in \mathcal{K}(\mathrm{T} \times \mathrm{T}')\)，函数 \(t' \mapsto f(t, t')\varphi(t, t')\) 是 \(\mu'\) -可积的（对几乎每个 \(t \in \mathrm{T}\)），且函数 \(t \mapsto \int f(t, t')\varphi(t, t') d\mu'(t')\) 几乎处处有定义，是 \(\mu\) -可积的；这里要用 Ch. III, §4, No.~1 的引理 1。其次注意，\(\varphi \mapsto \int d\mu(t) \int f(t, t')\varphi(t, t') d\mu'(t')\) 是 \(\mathrm{T} \times \mathrm{T}'\) 上以 \(\mu \otimes \mu'\) 为基的正测度，应用 Lebesgue-Nikodym 定理；最后利用 \(\mathcal{K}(\mathrm{T}')\) 中存在可数稠密集这一事实。）

\begin{enumerate}
\def\labelenumi{\alph{enumi})}
\setcounter{enumi}{1}
\item
  再设 T 可度量化。证明若满足：\(1^{\circ}\) 对几乎每个 \(t' \in T'\)，函数 \(t \mapsto f(t, t')\) 是 \(\mu\) -可测的；\(2^{\circ}\) 对几乎每个 \(t \in T\)，函数 \(t' \mapsto f(t, t')\)（关于 \(\mu'\)）几乎处处连续，则 a) 的条件成立。（利用 Ch. IV, §5 习题 13。）
\item
  取 \(\mathrm{T} = \mathrm{T}' = [0,1]\)。承认连续统假设（S, Ch. III, §6, No.~4），设 \(x \prec y\) 为 \(\mathrm{T}\) 上的一个良序关系，它没有最大元，且对每个 \(x \in \mathrm{T}\)，使 \(z \prec x\) 的 z 所成之集可数。若取 \(\mu = \mu'\) 为 Lebesgue 测度，证明特征函数 \(f\)（即偶 \((t, t')\) 中使 \(t \prec t'\) 的偶所成之集的特征函数）满足 \(a\))，但 \(\int f(t, t') d\mu(t) = 0\)（对每个 \(t' \in \mathrm{T}'\)），而 \(\int f(t, t') d\mu'(t') = 1\)（对每个 \(t \in \mathrm{T}\)）。
\end{enumerate}

\begin{enumerate}
\def\labelenumi{\arabic{enumi})}
\setcounter{enumi}{7}
\tightlist
\item
  设 \(\mathrm{T}\) 为局部紧空间，\(\mu\) 为一个 \(\neq 0\) 的正测度，定义在 \(\mathrm{T}\) 上，A 为离散空间，\(\lambda\) 为 \(\mathrm{A}\) 上以 \(+1\) 为 \(\mathrm{A}\) 中每点质量的测度。
\end{enumerate}

\begin{enumerate}
\def\labelenumi{\alph{enumi})}
\item
  为使映射 \(f\)（由 \(\mathrm{A} \times \mathrm{T}\) 到一个拓扑空间中）是 \((\lambda \otimes \mu)\) -可测的，必须且只需：对每个 \(\alpha \in \mathrm{A}\)，映射 \(t \mapsto f(\alpha, t)\) 是 \(\mu\) -可测的。
\item
  为使函数 \(\mathbf{f}\)（定义在 \(\mathrm{A} \times \mathrm{T}\) 上、取值于 \(\overline{\mathbf{R}}\) 或某个 Banach 空间中）是 \((\lambda \otimes \mu)\) -可积的，必须且只需：\(1^{\circ}\) 除去 \(\alpha \in \mathrm{A}\) 的一个可数值集外，函数 \(t \mapsto \mathbf{f}(\alpha, t)\) 恒等于零；\(2^{\circ}\) 对每个 \(\alpha \in \mathrm{A}\)，函数 \(t \mapsto \mathbf{f}(\alpha, t)\) 是 \(\mu\) -可积的；\(3^{\circ} \sum_{\alpha \in \mathrm{A}} \int |\mathbf{f}(\alpha, t)| d\mu(t) < +\infty\)。为使 \(\mathbf{f}\)
\end{enumerate}

是本质 \((\lambda \otimes \mu)\) -可积的，必须且只需：\(1^{\circ}\) 对每个 \(\alpha \in \mathbf{A}\)，函数 \(t \mapsto \mathbf{f}(\alpha, t)\) 是本质 \(\mu\) -可积的；\(2^{\circ} \sum_{\alpha \in \mathbf{A}} \int |\mathbf{f}(\alpha, t)| d\mu(t) < +\infty\)；此时

有

\[
\iint \mathbf {f} d \lambda d \mu = \sum_ {\alpha \in A} \int \mathbf {f} (\alpha , t) d \mu (t).
\]

\begin{enumerate}
\def\labelenumi{\alph{enumi})}
\setcounter{enumi}{2}
\tightlist
\item
  设 X 为局部紧空间，\((\alpha,t)\mapsto\rho_{\alpha,t}\) 为 X 上的一个正测度族 \((\alpha\in\mathrm{A},t\in\mathrm{T})\)。为使族 \((\alpha,t)\mapsto\rho_{\alpha,t}\) 是 \((\mu\otimes\nu)\) -适宜的，必须且只需：对每个 \(\alpha\in\mathrm{A}\)，族 \(t\mapsto\rho_{\alpha,t}\) 是 \(\mu\) -适宜的，且测度族 \(\left(\int\rho_{\alpha,t}d\mu(t)\right)_{\alpha\in\mathrm{A}}\) 是可和的。
\end{enumerate}

\begin{enumerate}
\def\labelenumi{\arabic{enumi})}
\setcounter{enumi}{8}
\tightlist
\item
  设 \(u\) 与 \(v\) 为 \(\mathbf{R}\) 上两个递增且右连续的数值函数，满足 \(u(x) = v(x) = 0\)（对 \(x < 0\)）。设 \(w\) 为 \(\mathbf{R}\) 上由下式定义的递增右连续函数：\(w(t) = u(t)v(t)\)（对 \(t \geqslant 0\)），\(w(t) = 0\)（对 \(t < 0\)）；设 \(\lambda, \mu, \nu\) 分别为与 \(u, v, w\) 相联系的 Stieltjes 测度（§6, 习题 5）。
\end{enumerate}

对每个函数 \(\mathbf{f}\)（定义在 \(\mathbf{R}\) 上、取值于 \(\overline{\mathbf{R}}\) 或某个 Banach 空间 \(\mathrm{F}\) 中），使它对应于函数 \(\overline{\mathbf{f}}\)（定义在 \(\mathbf{R}^2\) 上，由下述条件给出：\(\overline{\mathbf{f}}(x,y) = \mathbf{f}(x)\)（若 \(y < x\)），\(\overline{\mathbf{f}}(x,y) = \mathbf{f}(y)\)（若 \(y \geqslant x\)））。证明：为使 \(\mathbf{f}\) 是 \(\nu\) -可积的，必须且只需 \(\overline{\mathbf{f}}\) 关于乘积测度 \(\lambda \otimes \mu\) 可积，此时 \(\int \mathbf{f} d\nu = \iint \overline{\mathbf{f}} d\lambda d\mu\)（先对区间的特征函数证明）。由此推出公式

\[
\int \mathbf {f} (x) d w (x) = \int \mathbf {f} (x) v (x -) d u (x) + \int \mathbf {f} (x) u (x +) d v (x).
\]

特别地，若 \(u\) 与 \(v\) 在 \(\mathbf{R}\) 上连续，就得到分部积分公式

\[
\int_ {a} ^ {b} u (x) d v (x) = u (b) v (b) - u (a) v (a) - \int_ {a} ^ {b} v (x) d u (x).
\]

\begin{enumerate}
\def\labelenumi{\arabic{enumi})}
\setcounter{enumi}{9}
\tightlist
\item
  设 \(\mathbf{T}\) 与 \(\mathbf{X}\) 为两个局部紧空间，\(\lambda\) 为 \(\mathbf{T}\) 上的正测度，\(\mu\) 为 \(\mathbf{X}\) 上的有界正测度且 \(\mu(\mathbf{X}) = 1\)。设 \(f\) 为一个 \(> 0\) 的数值函数（在 \(\mathbf{T} \times \mathbf{X}\) 的每点），它与 \(\log f\) 都关于测度 \(\lambda \otimes \mu\) 可积。证明不等式
\end{enumerate}

\[
\log \left(\int \exp (\int \log f d \mu) d \lambda\right) \leqslant \int (\log \int f d \lambda) d \mu
\]

并证明：除非 \(f\) 等价于形如 \(g \otimes h\) 的函数，否则等式不可能成立（对每个 \(t \in \mathrm{T}\)，对函数 \(x \mapsto f(t,x) / (\int f(t,x)d\lambda (t))\) 应用几何平均不等式（Ch. IV, §6, 习题 7 d））。

\begin{enumerate}
\def\labelenumi{\arabic{enumi})}
\setcounter{enumi}{10}
\tightlist
\item
  设 \(p\) 为有限实数 \(\geqslant 1\)，\(T\) 与 \(X\) 为两个局部紧空间，\(\lambda\) 为 \(T\) 上的正测度，\(\mu\) 为 \(X\) 上的正测度，\(f\) 为函数 \(\geqslant 0\)（定义在 \(T \times X\) 上），它与 \(f^p\) 都关于测度 \(\lambda \otimes \mu\) 可积。证明不等式
\end{enumerate}

\[
\left(\int^ {*} \left(\int f d \mu\right) ^ {p} d \lambda\right) ^ {1 / p} \leqslant \int^ {*} \left(\int f ^ {p} d \lambda\right) ^ {1 / p} d \mu .
\]

（对每个 \(t \in \mathbf{T}\)，对写成下述形式的函数 \(x \mapsto f(t, x)\) 应用 Hölder 不等式

\[
f (t, x) = g (t, x) \left(\int f ^ {p} (t, x) d \lambda (t)\right) ^ {1 / p q},
\]

其中 \(q\) 是与 \(p\) 共轭的指数。）证明：除非 \(f\) 等价于形如 \(g \otimes h\) 的函数，否则等式不可能成立。

¶ 12) 设 \(\mathrm{T}_i\)（\(1 \leqslant i \leqslant n\)）为 \(n\) 个局部紧空间，\(\mu_i\) 为 \(\mathrm{T}_i\) 上的正测度（\(1 \leqslant i \leqslant n\)）。对每个指标 \(i\)，用 \(\mathrm{E}_i\) 表示乘积 \(\prod_{j \neq i} \mathrm{T}_j\)；设 \(f_i\)

为函数 \(\geqslant 0\)（在 \(\mathrm{T} = \prod_{i=1}^{n} \mathrm{T}_i\) 上），关于 \(\mu = \bigotimes_{i=1}^{n} \mu_i\) 可测且不依赖

于 \(t_{i}\)；证明：若对 \(1 \leqslant k \leqslant n\)，函数 \(f_{k}^{n-1}\) 关于测度 \(\mu_{1} \otimes \cdots \otimes \mu_{k-1} \otimes \mu_{k+1} \otimes \cdots \otimes \mu_{n}\) 可积，则 \(f_{1}f_{2} \cdots f_{n}\) 是 \(\mu\) -可积的，且

\[
\int f _ {1} f _ {2} \dots f _ {n} d \mu_ {1} d \mu_ {2} \dots d \mu_ {n} \leqslant \left(\prod_ {k = 1} ^ {n} \mathrm{J} _ {k}\right) ^ {1 / (n - 1)}
\]

其中，对每个指标 k，令 \(J_{k}=\int f_{k}^{n-1}d\mu_{1}\cdots d\mu_{k-1}d\mu_{k+1}\cdots d\mu_{n}\)（对 n 作归纳，应用 Lebesgue--Fubini 定理与 Hölder 不等式）。

由此推出：若 A 是 T 的 \(\mu\) -可测子集，\(A_{i}\) 是它在 \(E_{i}\) 上的投影，且 \(A_{i}\) 可积并具有测度 \(m_{i}\)（关于测度 \(\bigotimes_{j\neq i}\mu_{j}\)，在 \(E_{i}\) 上），则 A 是

\(\mu\) -可积的，且

\[
\mu (\mathrm{A}) \leqslant (m _ {1} m _ {2} \dots m _ {n}) ^ {1 / (n - 1)}.
\]

考察这两个不等式中取等号的情形。

推广到下述情形：不考察 \(n\) 个由 \(n - 1\) 个 \(\mathrm{T}_i\) 作成的乘积，而是考察 \(\binom{n}{p}\) 个由 \(p\) 个 \(\mathrm{T}_i\) 作成的乘积，并且在 \(\mathrm{T}\) 上积分 \(\binom{n}{p}\) 个函数 \(\geqslant 0\) 的乘积，其中每个函数只依赖 \(p\) 个变量 \(t_i\)。例如，若 \(\mathrm{F}_{ij} = \mathrm{T}_i \times \mathrm{T}_j (i < j)\) 且 \(f_{ij}\) 只依赖变量 \(t_i\) 与 \(t_j\)，证明

\[
\int \left(\prod_ {i <   j} f _ {i j}\right) d \mu_ {1} d \mu_ {2} \dots d \mu_ {n} \leqslant \left(\prod_ {i <   j} \int f _ {i j} ^ {n - 1} d \mu_ {i} d \mu_ {j}\right) ^ {1 / (n - 1)}.
\]

¶ 13) 设 \((T_{\iota})_{\iota \in I}\) 为一族紧空间，且对每个 \(\iota \in I\)，设 \(\mu_{\iota}\) 为 \(T_{\iota}\) 上总质量等于 1 的正测度。设 \(T\) 为乘积空间 \(\prod_{\iota \in I} T_{\iota}\)，

\(\mu\) 为 T 上的乘积测度 \(\bigotimes \mu_{\iota}\)（Ch. III, §4, No.~5）。

\begin{enumerate}
\def\labelenumi{\alph{enumi})}
\item
  对每个 \(\iota \in I\)，设 \(K_{\iota}\) 为 \(T_{\iota}\) 的紧子集；证明若 \(K = \prod_{\iota \in I} K_{\iota}\)，则 \(\mu(K) = \prod \mu_{\iota}(K_{\iota})\)。
\item
  设 M 为 T 的紧子集。证明存在 I 的可数子集 J，使得令 H = I - J，\(T_{J} = \prod_{\iota \in J} T_{\iota}\)，\(\mu_{J} = \bigotimes_{\iota \in J} \mu_{\iota}\)，\(T_{H} = \prod_{\iota \in H} T_{\iota}\) 时，有 \(M \subset N \times T_{H}\)，其中 N 是 \(\mu_{J}\) -可测的 \(T_{J}\) 子集，且 \(\mu_{\mathrm{J}}(\overline{\mathrm{N}}) = \mu(\mathrm{M})\)。（指出：对每个 n \textgreater{} 0，存在 M 的一个开邻域，它是有限个初等集的并，其测度与 \(\mu(\mathrm{M})\) 之差小于 1/n。）由此推出：对几乎每个 \(x \in N\)，截面 \(\mathrm{M}(x) \subset \mathrm{T}_{\mathrm{H}}\)（当把 T 与乘积 \(T_{J} \times T_{H}\) 等同时）包含诸测度 \(\mu_{\iota}\)（\(\iota \in H\)）的支集的乘积（利用 Lebesgue--Fubini 定理）。
\item
  证明：若 I 可数，且对每个 \(\iota \in I\)，\(A_{\iota}\) 是 \(\mu_{\iota}\) -可测的 \(T_{\iota}\) 子集，则集合 \(A = \prod_{\iota \in I} A_{\iota}\) 是 \(\mu\) -可测的，且 \(\mu(A) = \prod_{\iota \in I} \mu_{\iota}(A_{\iota})\)。
\item
  假设 I 不可数。对每个 \(\iota \in I\)，设 \(A_{\iota}\) 为 \(\mu_{\iota}\) -可测的 \(T_{\iota}\) 子集。为使 \(A = \prod_{\iota \in I} A_{\iota}\) 是 \(\mu\) -可测的，必须且只需处于
\end{enumerate}

下述两种情形之一：\(1^{\circ}\prod_{\iota \in I}\mu_{\iota}(A_{\iota}) = 0;2^{\circ}A_{\iota}\) 包含测度

\(\mu_{\iota}\) 的支集，但可能除去属于 I 的某个可数子集的指标 \(\iota\)。在这两种情形的每一种，\(\mu(A) = \prod_{\iota \in I} \mu_{\iota}(A_{\iota})\)。（假设不处于前述两种情形的任一种，

证明可化归于 \(\mu_{\iota}(\mathrm{A}_{\iota}) = 1\)（对所有 \(\iota \in \mathrm{I}\)）的情形。利用 \(b\)，证明 A 与 T-A 都不能包含具有非零测度的紧集。）

¶ 14) 设 K 为 R 中区间 {[}0,1{]}，λ 为 K 上的 Lebesgue 测度，A 为不可数集，μ 为 T = K\^{}A 上的测度，它是每个因子上的测度 λ 的乘积。

\begin{enumerate}
\def\labelenumi{\alph{enumi})}
\item
  证明 T 的每个闭的可度量化子空间 X 都是 \(\mu\) -可忽略的（利用上面的习题 13 b) 以及 GT, IX, §2 习题 8）。
\item
  设 Y 为具有可数基的局部紧空间，\(\nu\) 为 Y 上的正测度，\(\pi\) 为 Y 到 T 中的 \(\nu\) -逆紧映射。证明像 \(\pi(\nu)\) 与 \(\mu\) 互奇异（利用 a)，证明 \(\pi(\nu)\) 集中于一个 \(\mu\) -可忽略集上）。
\end{enumerate}

\begin{enumerate}
\def\labelenumi{\arabic{enumi})}
\setcounter{enumi}{14}
\tightlist
\item
  设 \((\mathrm{T}_{\iota})_{\iota \in I}\) 为任意一族紧空间，且对每个 \(\iota \in I\)，设 \(\mu_{\iota}\) 为 \(\mathrm{T}_{\iota}\) 上总质量为 1 的正测度；设 \(\mu\) 为乘积测度 \(\bigotimes_{\iota \in I} \mu_{\iota}\)（在 \(\mathrm{T} = \prod_{\iota \in \mathrm{I}}\mathrm{T}_{\iota}\) 上）。对把 I 分成两个子集的每个分割 (L, M)，把 T 与
\end{enumerate}

乘积 \(\mathrm{T_L} \times \mathrm{T_M}\) 等同，其中 \(\mathrm{T_L} = \prod_{\iota \in \mathrm{L}} \mathrm{T}_{\iota}\)；对每个 \(t \in \mathrm{T}\)，记 \(t_{\mathrm{L}} = \mathrm{pr}_{\mathrm{L}} t\)，使得

\(t\) 与 \((t_{\mathrm{L}}, t_{\mathrm{M}})\) 等同；设 \(\mu_{\mathrm{L}}\) 为测度 \(\bigotimes_{\iota \in \mathrm{L}} \mu_{\iota}\)（在 \(\mathrm{T}_{\mathrm{L}}\) 上）。设 \(p\) 为有限

实数 \(\geqslant 1\)；对每个函数 \(\mathbf{f}\)（属于 \(\mathcal{L}_{\mathrm{F}}^{p}(\mathrm{T},\mu)\)；F 为 Banach 空间或 \(\mathrm{F} = \overline{\mathbf{R}}\)），记 \(\mathbf{f}_{\mathrm{L}}(t) = \int \mathbf{f}(t_{\mathrm{L}},t_{\mathrm{M}}) d\mu_{\mathrm{M}}(t_{\mathrm{M}})\)。证明：关于 I 的有限子集 L 所成的有向集，\(\mathbf{f}_{\mathrm{L}}\) \(p\) 阶平均收敛于 \(\mathbf{f}\)，而 \(\mathbf{f}_{\mathrm{M}}\) \(p\) 阶平均收敛于等于 \(\int \mathbf{f} d\mu\) 的常值函数。（用一个只依赖有限个变量的连续函数逼近 \(\mathbf{f}\)。）

由此推出：若 \(\mu\) -可测子集 \(A\)（含于 \(T\)）满足：对一切 \(t \in A\)，每个点 \(t' \in T\)（其坐标除有限个指标外都与 \(t\) 的坐标相等）也属于 \(A\)，则 \(\mu(A) = 0\) 或 \(\mu(A) = 1\)。

¶ 16) 设 \((T_n)\) 为紧空间的无穷序列，\(\mu_n\) 为 \(T_n\) 上总质量等于 1 的正测度，\(\mu\) 为乘积测度 \(\bigotimes_{n=1}^{\infty} \mu_n\)（在 \(T = \prod_{n=1}^{\infty} T_n\) 上）。

\begin{enumerate}
\def\labelenumi{\alph{enumi})}
\item
  设 \(f\) 为函数 \(\geqslant 0\)，\(\mu\) -可积（在 \(\mathbf{T}\) 上）。设 \((\mathrm{L}_n)\) 为 \(\mathbf{N}\) 的子集的递增序列，且令 \(\mathrm{M}_n = \mathbf{N} - \mathrm{L}_n\)，设 \(g = \sup_n f_{\mathrm{L}_n}\)，\(h = \sup_n f_{\mathrm{M}_n}\)（记号见习题 15）。对每个 \(\alpha > 0\)，设 \(\mathrm{A}_{\alpha}\) 为 \(t \in \mathbb{T}\) 中使 \(g(t) > \alpha\) 的点所成之集，\(\mathrm{B}_{\alpha}\) 为 \(t \in \mathbb{T}\) 中使 \(h(t) > \alpha\) 的点所成之集。证明 \(\alpha \cdot \mu(\mathrm{A}_{\alpha}) \leqslant \int f d\mu\) 且 \(\alpha \cdot \mu(\mathrm{B}_{\alpha}) \leqslant \int f d\mu\)。（指出：\(\mathrm{A}_{\alpha}\) 是 \(t \in \mathbb{T}\) 中使诸 \(f_{\mathrm{L}_n}(t)\) 里至少有一个 \(> \alpha\) 的点所成之集，把它表为两两不交的集 \(\mathrm{G}_n\) 的可数并，使 \(\alpha \cdot \mu(\mathrm{G}_n) \leqslant \int_{\mathrm{G}_n} f d\mu\)。）
\item
  假设 \((\mathrm{L}_n)\) 是 \(\mathbf{N}\) 的有限子集的递增序列，其并为 \(\mathbf{N}\)。证明 \(f_{\mathrm{L}_n}\) 几乎处处收敛于 \(f\)，且 \(f_{\mathrm{M}_n}\) 在 \(\mathrm{T}\) 中几乎处处收敛于常数 \(\int f d\mu\)。（对每个 \(\varepsilon > 0\)，考察一个连续函数 \(g\)（只依赖有限个变量且满足 \(\int |f - g| d\mu \leqslant \varepsilon\)），并把 \(a\) 应用于函数 \(|f - g|\)。）
\end{enumerate}

\begin{enumerate}
\def\labelenumi{\arabic{enumi})}
\setcounter{enumi}{16}
\tightlist
\item
  设 \((\mathrm{T}_{\iota})_{\iota \in \mathrm{I}}\) 为任意一族紧空间；对每个 \(\iota \in \mathrm{I}\)，设 \(\mu_{\iota}\) 为 \(\mathrm{T}_{\iota}\) 上总质量等于 1 的正测度，且设 \(\mu\) 为乘积测度 \(\bigotimes_{\iota \in \mathrm{I}} \mu_{\iota}\)（在 \(\mathrm{T} = \prod_{\iota \in \mathrm{I}} \mathrm{T}_{\iota}\) 上）。对每个 \(\iota \in \mathrm{I}\)，设 \(f_{\iota}\) 为函数 \(\geqslant 0\)（定义在 \(\mathrm{T}_{\iota}\) 上），\(\mu_{\iota}\) -可积且满足 \(\int f_{\iota} d\mu_{\iota} \leqslant 1\)；令 \(\mu_{\iota}' = f_{\iota} \cdot \mu_{\iota}\) 且 \(\mu' = \bigotimes_{\iota \in \mathrm{I}} \mu_{\iota}'\)（Ch. III, §4, 习题 6）。假设 \(\mu' \neq 0\)。
\end{enumerate}

\begin{enumerate}
\def\labelenumi{\alph{enumi})}
\tightlist
\item
  对每个有限子集 \(J\)（含于 \(I\)），令 \(T_J = \prod_{\iota \in J} T_\iota\)，\(\mu_J = \bigotimes_{\iota \in J} \mu_\iota\)，\(\mu_J' = \bigotimes_{\iota \in J} \mu_\iota'\)，
\end{enumerate}

\(f_{\mathrm{J}}(t) = \prod_{\iota \in \mathbb{J}}f_{\iota}(\mathrm{pr}_{\iota}t), g_{\mathrm{J}} = \sqrt{f_{\mathrm{J}}}, \mu_{\mathrm{J}}^{\prime \prime} = \sqrt{\mu_{\mathrm{J}}\mu_{\mathrm{J}}^{\prime}}\)（§5, No.~9），且 \(\rho (\mu_{\mathrm{J}},\mu_{\mathrm{J}}^{\prime}) = \mu_{\mathrm{J}}^{\prime \prime}(\mathrm{T}_{\mathrm{J}}) =\)

\(\int g_{\mathrm{J}}d\mu_{\mathrm{J}}\)。证明：为使诸函数 \(g_{\mathrm{J}}\) 关于 I 的有限子集所成的有向有序集均方收敛于 \(\mathcal{L}^2 (\mathrm{T},\mu)\) 中的一个函数，必须且只需数 \(\rho (\mu_{\iota},\mu_{\iota}^{\prime})\) 的乘积在 \(\mathbf{R}_{+}^{*}\) 中收敛，即 \(>0\)。（对满足 \(\mathrm{J}\subset \mathrm{L}\) 的两个有限子集 J, L，估计范数 \(\mathrm{N}_2(g_{\mathrm{J}} - g_{\mathrm{L}})\)（借助于诸 \(\rho (\mu_{\iota},\mu_{\iota}^{\prime})\)）。）由此推出：诸函数 \(f_{\mathrm{J}}\) 此时平均收敛于一个函数 \(f\geqslant 0\)，它满足 \(\mu^{\prime} = f\cdot \mu\)。

\begin{enumerate}
\def\labelenumi{\alph{enumi})}
\setcounter{enumi}{1}
\tightlist
\item
  证明：若乘积 \(\prod_{\iota \in I} \rho(\mu_{\iota}, \mu_{\iota}')\) 为零，则测度 \(\mu\) 与 \(\mu'\) 互
\end{enumerate}

奇异（考察集 \(\mathrm{A_J}\)（\(t\in \mathrm{T}\) 中使 \(g_{\mathrm{J}}(t) > 1\) 的点所成之集），并估计测度 \(\mu (\mathrm{A_J})\) 与 \(\mu '(T - A_J)\)）。

¶ 18) 记号与习题 17 相同，设 \(\nu_{\iota}\) 为 \(\mathrm{T}_{\iota}\) 上总质量等于 1 的正测度，且设 \(\nu = \bigotimes \nu_{\iota}\)。

\begin{enumerate}
\def\labelenumi{\alph{enumi})}
\item
  证明：若诸测度 \(\nu_{\iota}\) 中有一个与具有同一指标的测度 \(\mu_{\iota}\) 互奇异，则 \(\nu\) 与 \(\mu\) 互奇异。
\item
  对每个 \(\iota \in I\)，写 \(\nu_{\iota} = \mu_{\iota}' + \mu_{\iota}''\)，其中 \(\mu_{\iota}'\) 以 \(\mu_{\iota}\) 为基，而 \(\mu_{\iota}''\) 与 \(\mu_{\iota}\) 互奇异（§5, 定理 3）；假设 \(\mu_{\iota}' \neq 0\)（对所有 \(\iota \in I\)）。证明：为使 \(\nu\) 不与 \(\mu\) 互奇异，必须且只需 \(\mu_{\iota}'' = 0\) 除去一个可数指标族外成立，且乘积 \(\mu' = \bigotimes_{\iota \in I} \mu_{\iota}'\) 是一个测度 \(\neq 0\)，以 \(\mu\) 为基；此时 \(\mu'' = \nu - \mu'\)
\end{enumerate}

与 \(\mu\) 互奇异。（若对不可数无穷多个指标有 \(\mu_{\iota}^{\prime\prime} \neq 0\)，证明存在数 \(\alpha\) 使 \(0 < \alpha < 1\)，且对可数无穷多个指标 \(\iota_{n}\) 有 \(\mu_{\iota_{n}}^{\prime}(T_{\iota_{n}}) \leqslant \alpha\)；然后证明 \(\mu\) 与 \(\nu\) 互奇异（利用 \(a\)）；命题的第二部分利用习题 17 证明。）

\begin{enumerate}
\def\labelenumi{\arabic{enumi})}
\setcounter{enumi}{18}
\tightlist
\item
  设 \(X\) 为局部紧空间，\(Y\) 为无穷远处可数的局部紧空间，\(A\) 为 \(X \times Y\) 的普遍可测子集。
\end{enumerate}

\begin{enumerate}
\def\labelenumi{\alph{enumi})}
\item
  对每个 \(x \in X\)，证明截面 \(\mathrm{A}(x)\) 是 Y 的普遍可测子集。此外，对 Y 上的每个正测度 \(\mu\)，函数 \(x \mapsto \mu^{*}(\mathrm{A}(x))\) 在 X 上普遍可测（利用 Lebesgue--Fubini 定理）。
\item
  设 \(\mu\) 为 Y 上的正测度，使得对几乎每个 \(y \in Y\)（关于 \(\mu\)），A 的截面 \(\overline{A}(y)\) 可数。证明 \(x \in X\) 中使 \(\mu^{*}(A(x)) > 0\) 的点所成之集 N 不能包含任何可数紧集。
\end{enumerate}

\begin{enumerate}
\def\labelenumi{\arabic{enumi})}
\setcounter{enumi}{19}
\tightlist
\item
  \begin{enumerate}
  \def\labelenumii{\alph{enumii})}
  \tightlist
  \item
    设 \(\nu\) 为单位圆 \(\mathbf{U}: |z| = 1\)（在 \(\mathbf{C}\) 中）上的正测度，\(\mu\) 为复测度 \(j \cdot \nu\)，其中 \(j: z \mapsto z\) 是 \(\mathbf{U}\) 到 \(\mathbf{C}\) 中的典范单射，于是 \(|\mu| = \nu\)。设 \(\rho\) 为 \(\nu\) 在典范局部同胚 \(\varphi: t \mapsto e^{it}\)（由 \(\mathbf{R}\) 到 \(\mathbf{U}\) 上）下的逆像（§6, No.~6）。设 A 为 \(\nu\) -可测的 \(\mathbf{U}\) 子集；若 \(\mu(A) = |\mu(A)|e^{i\omega}\)，证明 \(|\mu(A)| \leqslant \left|\int_{\omega - \pi/2}^{\omega + \pi/2} e^{i\theta} d\rho(\theta)\right|\)。
  \end{enumerate}
\end{enumerate}

\begin{enumerate}
\def\labelenumi{\alph{enumi})}
\setcounter{enumi}{1}
\tightlist
\item
  由 \(a\) 推出：当 A 跑遍 \(\mu\) -可测的 \(\mathbf{U}\) 子集所成之集时，有
\end{enumerate}

\[
\begin{array}{l} \sup _ {\mathrm{A}} | \mu (\mathrm{A}) | \geqslant \sup _ {\omega} \left| \int_ {\omega - \pi / 2} ^ {\omega + \pi / 2} \mathcal {R} \big (e ^ {i (\theta - \omega)} \big) d \rho (\theta) \right| \\ \geqslant \frac {1}{2 \pi} \int_ {0} ^ {2 \pi} d \omega \int_ {\omega - \pi / 2} ^ {\omega + \pi / 2} \mathcal {R} \big (e ^ {i (\theta - \omega)} \big) d \rho (\theta) \end{array}
\]

并由此外推出

\[
\sup _ {A} | \mu (A) | \geqslant \frac {1}{\pi} \| \mu \|\tag{*}
\]

（借助 Lebesgue--Fubini 定理计算最后一个积分）。

\begin{enumerate}
\def\labelenumi{\alph{enumi})}
\setcounter{enumi}{2}
\tightlist
\item
  证明不等式 (*) 对局部紧空间 X 上的每个有界复测度都成立，其中 A 跑遍 X 的 \(\mu\) -可测子集所成之集。（设 \(\mu = h \cdot |\mu|\)，其中 h 取值于 U。考察 U 上的像测度 \(\nu = h(|\mu|)\)，并证明 \(h(\mu) = j \cdot \nu\)；然后对 \(h(\mu)\) 应用 b)。）
\end{enumerate}

\begin{enumerate}
\def\labelenumi{\arabic{enumi})}
\setcounter{enumi}{20}
\tightlist
\item
  用 \(\lambda\) 表示 \(\mathbf{R}^n\) 上的 Lebesgue 测度。
\end{enumerate}

\begin{enumerate}
\def\labelenumi{\alph{enumi})}
\tightlist
\item
  设 \(\mathbf{C}_0\) 为 \(\mathbf{R}^n\) 中的一个立方体（GT, VI, §1, No.~1），\(u\) 为 \(\mathbf{C}_0\) 上（关于 \(\lambda\)）可积的数值函数。设 \(s\) 为满足下式的实数
\end{enumerate}

\[
s \geqslant \frac {1}{\lambda (\mathrm{C} _ {0})} \int_ {\mathrm{C} _ {0}} | u | d \lambda .
\]

证明存在开立方体的一个序列 \((\mathrm{I}_k)\)（含于 \(\mathbf{C}_0\)），两两不交，且满足以下条件：

\(1^{\circ}\left|u(x)\right|\leqslant s\) 在集合 \(\mathrm{C_0 - }\bigcup_{k}\mathrm{I}_k;\) 中几乎处处成立；

\(2^{\circ}\) 若令

\[
u _ {k} = \frac {1}{\lambda (\mathrm{I} _ {k})} \int_ {\mathrm{I} _ {k}} | u (x) | d \lambda (x),
\]

则 \(u_{k} \leqslant 2^{n}s\) 对所有 \(k\) 成立；

\[
3 ^ {\circ} \sum_ {k} \lambda (I _ {k}) \leqslant \frac {1}{s} \int_ {C _ {0}} | u (x) | d \lambda (x).
\]

（把 \(C_{0}\) 的每个投影分成两个相等的区间，在 \(2^{n}\) 个两两不交的开立方体中选出立方体 \(I_{1i}\)，使相应的“平均值”\(u_{1i}\) 都 \(\geqslant s\)，并利用 s 的取法证明 \(u_{1i} \leqslant 2^{n}s\)。类似地，把第一次分割中每个立方体分解成 \(2^{n}\) 个立方体，但诸 \(I_{1i}\) 中的立方体除外，这样得到第二个有限的立方体族 \(I_{2i}\)，它们是在所得到的全部新立方体中选出的、使相应平均值 \(u_{2i}\) \(\leqslant s\) 的那些立方体。如此递归地继续下去。为验证条件 \(1^{\circ}\)，用反证法，注意子集 \(B_{m}\)——它由 \(C_{0} - \bigcup_{k} I_{k}\) 中使 \(|u(x)| \geqslant s + 1/m\) 的点组成——除去一个可忽略集外，是开集的递减序列 \((\mathrm{U}_r)\) 之交，其中 \(\mathrm{U}_r\) 是第 \(r\) 次分割中不属于诸 \(\mathrm{I}_k\) 且包含 \(\mathrm{B}_m\) 中一点的那些立方体之并。

\begin{enumerate}
\def\labelenumi{\alph{enumi})}
\setcounter{enumi}{1}
\tightlist
\item
  设 \(\mathcal{F}\) 为可积数值函数 \(u\)（定义在 \(C_0\) 上）所成的、具有下述性质的集：对每个立方体 \(C \subset C_0\)，若 \(m_C(u)\) 是平均值
\end{enumerate}

\[
\frac {1}{\lambda (\mathrm{C})} \int_ {\mathrm{C}} u (x) d \lambda (x),
\]

则 \(\int_{\mathbb{C}} |u(x) - m_{\mathbb{C}}(u)| d\lambda(x) \leqslant \lambda(\mathbb{C})\)；若 \(u \in \mathcal{F}\)，则 \(u - c\)（对每个实数 \(c\)）亦然。对每个函数 \(u \in \mathcal{F}\)、每个数 \(\sigma > 0\) 与每个立方体 \(\mathbb{C} \subset \mathbb{C}_0\)，用 \(\mathrm{S}(\sigma, u, \mathbb{C})\) 表示 \(x \in \mathbb{C}\) 中使 \(|u(x) - m_{\mathbb{C}}(u)| \geqslant \sigma\) 的点所成之集。最后，用 \(\mathrm{G}(\sigma)\) 表示使

\[
\lambda \bigl (\mathrm{S} (\sigma , u, \mathrm{C}) \bigr) \leqslant \mathrm{G} (\sigma) \int_ {\mathrm{C}} | u (x) - m _ {\mathrm{C}} (u) | d \lambda (x)
\]

（当 \(u\) 跑遍 \(\mathcal{F}\)、\(C\) 跑遍含于 \(C_0\) 的立方体全体时）成立的最小数；有 \(G(\sigma) \leqslant 1 / \sigma\)。证明：若 \(\sigma / 2^n > s \geqslant 1\)，则

\[
\mathrm{G} (\sigma) \leqslant \frac {1}{s} \mathrm{G} (\sigma - 2 ^ {n} s).
\]

（只需证明：对某个 \(u \in \mathcal{F}\)，有

\[
\lambda \left(\mathrm{S} (\sigma , u, \mathrm{C} _ {0})\right) \leqslant \frac {1}{s} \mathrm{G} (\sigma - 2 ^ {n} s) \int_ {\mathrm{C} _ {0}} | u (x) - m _ {\mathrm{C} _ {0}} (u) | d \lambda (x),
\]

并且为简单起见可设 \(m_{\mathrm{C}_0}(u) = 0\)。然后利用把 \(\mathrm{C}_0\) 分解为诸立方体 \(\mathrm{I}_k\) 之并以及一个使 \(|u(x)| \leqslant s\) 几乎处处成立的集合的分解，注意：若 \(x \in \mathrm{I}_k\) 属于 \(\mathrm{S}(\sigma, u, \mathrm{C}_0)\)，则 \(|u(x) - m_{\mathrm{I}_k}(u)| \geqslant \sigma - 2^n s\)，从而

\[
\lambda \left(\mathrm{I} _ {k} \cap \mathrm{S} (\sigma , u, \mathrm{C} _ {0})\right) \leqslant \mathrm{G} (\sigma - 2 ^ {n} s) \int_ {\mathrm{I} _ {k}} | u (x) - m _ {\mathrm{I} _ {k}} (u) | d \lambda (x). \Big)
\]

\begin{enumerate}
\def\labelenumi{\alph{enumi})}
\setcounter{enumi}{2}
\tightlist
\item
  由 \(b\) 推出：存在两个常数 \(B, b\)，只依赖于 \(n\)，使得对每个函数 \(u \in \mathcal{F}\) 与每个数 \(\sigma > 0\)，有 \(\lambda(\mathrm{S}(\sigma, u, \mathrm{C}_0)) \leqslant \mathrm{Be}^{-b\sigma}\lambda(\mathrm{C}_0)\)。（取 \(s = e\)（在 \(b\) 中）。）
\end{enumerate}

